\section{The Life of Reason}

\paragraph{Succession}
\textbf{Bishop Kallistos Ware} points out:

\begin{quotex}
Alongside this apostolic succession, largely hidden, existing on a `charismatic' rather than an official level, there is secondly the apostolic succession of the spiritual fathers and mothers in each generation of the Church—the succession of the saints, stretching from the apostolic age to our own day, which Saint Symeon the New Theologian termed the ``golden chain."… Both types of succession are essential for the true functioning of the Body of Christ, and it is through their interaction that the life of the Church on earth is accomplished. 

\end{quotex}
The West likewise has an esoteric hierarchy in addition to the exoteric hierarchy, although the history of the distinction is less precise. The charismatic succession has sometimes been secret, and then visible, sometimes in opposition to the exoteric Church, other times in harmony. \textbf{Valentin Tomberg} identified this as the \emph{Hermetic Tradition}, a parallel succession to the visible church, but not in opposition to it.

\textbf{Vladimir Soloviev} recognized the hidden charismatic succession, seeing himself as the heir of Kabbalah, Neoplatonism, Boehme, Swedenborg, and the German philosopher Friedrich Schelling. \textbf{Wolfgang Smith} followed a similar chain, but replaced Swedenborg with Meister Eckhart.

Yet, as we have seen, it goes back further, even to the pagans by way of the Zoroastrian Magi. Solovyov pointed out that the primary doctrines of Christianity — the Trinity and the Incarnation of the Logos — were known to the ancients. Hence, what Christianity added was to make those doctrines historical, concreate realities. That makes all the difference.

In recent times, Solovyov, Tomberg, Mouravieff, and Wolfgang Smith, \emph{inter alios}, are all laymen who have responded to the Spirit to keep the esoteric succession a living force. The Spirit blows where He wills.

\paragraph{The Life of Reason}
A common assumption is that the spiritual life is somehow irrational, relying on nothing but blind faith or arbitrary belief. Nothing could be further from the truth.

\textbf{David Hume}, the skeptic and alleged rationalist thinker, asserted that ``reason is and ought to be the slave of the passions". That is, a slave to whatever arises by chance. Essentially, Hume claims that man acts, and should act, like a highly intelligent dog. This is subrational. For Hume, the function of the intellect is to choose the best means to reach an end determined by the instinctive life.

Many call this ``living in the moment", ``life affirming", and so on, all delusions to rationalize the need to remain a slave. Living like this is actually contrary to nature, since it is man's nature to be intelligent, not instinctive, to be free, not a slave.

The spiritual life, on the other hand, is supra-rational. The intellect must be trained to dominate the animal and vegetable souls, i.e., one's animal and instinctive life. Its nature is to rule over the passions. Specifically, the intellect chooses the ends, not just the means. This is living in harmony with nature.

\paragraph{Divine Will}
The Divine Will is the source of the cosmic order and is one with the Logos, or the Principle of cosmic order. So the Life of Reason — in which a person chooses not just the means, but also the end — is to live in harmony with that cosmic order.

No one can elude this Will, whatever he may think, so everyone is in the place suitable for him. In the \emph{Symbolism of the Cross}, \textbf{Rene Guenon} points out that the Universal Will divides human beings into three categories, dependent on their reaction to the cosmic order.

\begin{itemize}
\item \textbf{The Faithful} conform consciously and voluntarily to the Universal Will. 
\item \textbf{The Infidels} obey the Law only against their will. 
\item \textbf{The Ignorant} have no clue. 
\end{itemize}
Unfortunately, in our age, our ``real" place in the cosmic order seldom matches our nominal place. Hence, it takes considerable intellectual effort to determine the real nature of things. We can no longer judge by outward appearances. Rene Guenon makes this clear in the \emph{Crisis of the Modern World}.

\begin{quotex}
Have we not arrived at that terrible age, announced in the Sacred Books of India, ``when the castes shall be mingled, when even the family shall no longer exist"?

As we have already pointed out, under the present state of affairs in the Western world, nobody any longer occupies the place that he should normally occupy by virtue of his own nature; this is what is meant by saying that the castes no longer exist, for caste, in its traditional meaning, is nothing other than individual nature, with the whole array of special aptitudes that this carries with it and that predisposes each man to the fulfillment of one or another particular function. 

\end{quotex}
\paragraph{Rhetoric}
\begin{quotex}
An argument based on knowledge implies instruction, and there are people whom one cannot instruct. \flright{\textsc{Aristotle}}

\end{quotex}
Dialectic and Rhetoric are two traditional sciences of the Trivium (along with Grammar). Dialectic is concerned with Logic and the mechanics of thought. Rhetoric is about the means of persuasion.

Imagine a possible world in which journalists were skilled in Dialectic and Rhetoric. Then, in their political commentaries, they could point out the logical flaws and rhetorical devices in the speeches of politicians. Although journalists are seldom logical, they do, to some extent, understand rhetoric. Hence, they will judge whether rhetoric has been effective or not in its aim of persuasion.

However, we, as the ``informed public" need to be the impartial judges. Dialectic deals with logic and rhetoric with facts. Unfortunately, Aristotle points out that rhetoric is little more than:

\begin{quotex}
The arousing of prejudice, pity, anger, and similar emotions [which] has nothing to do with the essential facts, but is merely a personal appeal to the man who is judging the case. 

\end{quotex}
It is necessary to ignore such appeals, for the only things of importance are whether

\begin{itemize}
\item Something has happened or has not happened, 
\item What will be or will not be. 
\end{itemize}
Thus, in a contemporary political context, the primary concern should be with results, past, present, and future. The human race lives metaphysically far from the Center of the Universe. Hence, we are subjected to laws at several different levels: the natural law, economic laws, social laws, psychological laws, laws of biology, chemistry and physics. This must all be taken into consideration in judging what may or may not happen.

Ideally, then, in the Life of Reason, Rhetoric must be in the service of knowledge. Since people have different abilities and intelligence levels, rhetoric needs to be adapted to a particular audience. The audience will be swayed or moved by rhetoric to support a position or course of action. Of course, the expert is neither swayed nor moved, but is only concerned about the actual effects of that course. The speaker will use enthymemes, i.e., unspoken parts of a syllogism that the audience is expected to fill in. Enthymemes make the audience members feel ``smart", as though they were in on special knowledge. The intelligent listener will be on the alert for such enthymemes.

\flrightit{Posted on 2016-09-15 by Cologero }
